% Shared book design. Content and ordering live in Markdown and src/book.json.
\documentclass[11pt,twoside,openany]{scrbook}
\usepackage[
  mode=book,lang=zh,theme=jade,background=syn-paper,
  font-profile=pagella,cjk-profile=source-han,
  toc=false,title-layout=page,numbering=none,theorem-style=boxed,
  density=balanced,measure=wide,card-radius=3pt,rail-width=2pt,rule-weight=0.4pt
]{synaptic}
\setcounter{tocdepth}{0}
\emergencystretch=2em
\allowdisplaybreaks[2]
\SynapticTitle{全国大学生数学竞赛题解}
\SynapticHeader{试题\quad ·\quad 解答\quad ·\quad 评注}
\SynapticShortTitle{数学竞赛题解}
\SynapticSubtitle{初赛与决赛\quad ·\quad 思想、方法与结构}
\SynapticAuthor{Jacobi Lie}
\SynapticEmail{jssnlzx@seu.edu.cn}
\SynapticDate{2026}
\SynapticSubject{数学竞赛题解与方法}
\SynapticKeywords{数学竞赛,分析,代数,几何}
% One title per problem, with a stable PDF destination independent of display order.
\makeatletter
\newcommand{\BookProblem}[2]{%
  \par\addvspace{1.2em}\Needspace{7\baselineskip}%
  \phantomsection
  \protected@edef\@currentlabel{#1}%
  \def\@currentlabelname{#1}%
  \label{pr:#2}%
  \hypertarget{pr:#2}{}%
  \pdfbookmark[1]{#1}{problem.#2}%
  \markright{#1}%
}
\makeatother
\newcommand{\BookProblemEnd}{\par}
